\textbf{Figure 6: Lung-cancer case study --- differential metabolomics by mutation status and OmniPath Metabo Azelate query.}
A worked use case applying the OmniPath Metabo database to the Shorthouse \emph{et al.}\ 2022 lung-cancer cohort (CCLE proteomics + targeted metabolomics) demonstrates how the integrated database supports differential-metabolomics interpretation across mutation contrasts and downstream evidence retrieval.
\textbf{(a)} KRAS-mutated vs. wild-type volcano plot: log\textsubscript{2} fold-change against $-\log_{10}$ $P$-value for every measured metabolite feature in the KRAS contrast (limma). Features passing $|\log_2 \text{FC}| > 0.5$ and $P < 0.05$ are highlighted; remaining features are shown in light grey.
\textbf{(b)} EGFR-mutated vs. wild-type volcano on identical axis limits as (a) so the two contrasts are visually comparable (FR-019).
\textbf{(c)} OmniPath Metabo evidence for Azelate, a top EGFR differential metabolite: relation counts by interaction type, stacked by upstream resource (ChEMBL, DrugCentral, FooDB, HMDB, MACDB, MetAtlas, PFOCR, RHEA, STITCH, ChEBI, SwissLipids). Data are restricted to the resolved canonical \emph{Azelate} entity (InChI key \texttt{BDJRBEYXGGNYIS-UHFFFAOYSA-N}); the two source-specific seed records are retained in the supplementary tables (\textit{Azelaic acid (PTF52203)}, \textit{Azelaic acid; nonanedioic acid}) but excluded from the plotted counts.
\textbf{(d)} Cancer-association evidence for the resolved \emph{Azelate} entity: evidence-record counts per disease type stacked by sample (tissue) type (Blood, Serum, Plasma, Tissue, Stool, Urine, Sweat, Cells). Counts are direct MACDB extractions through the OmniPath Metabo FastAPI surface.
